\section{Problem Setting and Pipeline Overview}
\label{sec:setting}

\subsection{Notation}
\label{sec:notation}

A capture session provides $V$ calibrated views.  View $v$ has camera $\cam_v$ (intrinsics
$f_x, f_y$, principal point, world-to-camera rotation $R_v$ and translation), image domain
$\Omega_v = [0, W_v) \times [0, H_v)$, and --- during capture only --- a photograph
$I_v : \Omega_v \to \R^3$.  Stage~1 fits each photograph with a compact 2D Gaussian field
$\field_v$ (\cref{sec:compact-capture}); afterwards the photographs are discarded.  The
reconstruction state is a set of $N$ 3D Gaussians
$\{(\bm{\mu}_i, \Sigma_i, o_i, \bm{c}_i)\}_{i=1}^{N}$ with means
$\bm{\mu}_i \in \R^3$, covariances $\Sigma_i \in \Sym(3)$, opacities $o_i \in (0,1)$, and
SH color coefficients $\bm{c}_i$, rendered exactly as in standard 3DGS.

\begin{definition}[Reconstruction inputs]
\label{def:recon-inputs}
$\mathcal{R} = \{(\cam_v, \field_v)\}_{v=1}^{V}$: calibrated cameras plus frozen per-view 2D
Gaussian observation fields.  No RGB image, alpha mask, or dense tensor is part of
$\mathcal{R}$.
\end{definition}

The data boundary is structural, not aspirational: the reconstruction entry points accept only
$\mathcal{R}$ and configuration; their implementations do not import an image loader or the
dense RGB trainer, do not open image suffixes, and reject dense handovers.  Tests deny these
imports at runtime, so a violation is a test failure, not a code-review judgment
call.\evidence{ara/logic/claims.md C31; src/rtgs/carrier_pipeline.py; tests/test_pipeline.py}

\subsection{The two arms and the RGB control}
\label{sec:arms}

\begin{figure}[t]
  \centering
  \figplaceholder{7.5cm}{%
    \textbf{Figure 1 --- pipeline overview (teaser).}  Three-lane diagram.  Lane 1 (capture,
    grayed out to mark it as upstream of the claim boundary): source RGB $\to$ per-view 2D
    Gaussian fit $\to$ frozen compact captures; RGB shown being discarded (crossed out).
    Lane 2 (Arm 1, no Beam): captures $\to$ Beam-independent 3D initialization (balanced
    compact carve) $\to$ compact-field-supervised 3DGS $\to$ final 3D Gaussians + memory
    receipt.  Lane 3 (Arm 2, Beam): the same captures $\to$ Beam Fusion $\to$ carrier
    refinement $\to$ the same compact optimizer.  A fourth thin lane shows the RGB control
    (standard image-supervised 3DGS from the identical frozen initialization).  Annotate the
    claim boundary (``post-fit reconstruction VRAM'') with a vertical dashed line after the
    fit.  Include per-lane byte sizes for one real scene once measured
    (\cref{sec:vram-protocol}).  Suggested source: draw as TikZ or vector graphics; render
    example thumbnails from \repopath{dataset/} scene frame\_00008 captures.}
  \caption{\redcaption{Pipeline overview and claim boundary.  Figure to be produced; the
    layout is specified in the placeholder.}}
  \label{fig:teaser}
\end{figure}

Both reconstruction arms consume the identical $\mathcal{R}$ and differ only in
initialization; the RGB control isolates the supervision variable.

\textbf{Arm 1 --- direct compact reconstruction (no Beam).}
Balanced compact-carve initialization with a hard 3D count cap
(\cref{sec:exp1-baselines}), then sampled compact-field optimization
(\cref{sec:compact-supervision}) under a frozen no-Beam schedule, then compact-domain
validation and standard \mbox{3DGS} output.  This arm must never import or call Beam Fusion,
Beam contributor lineage, carrier covariance repair, or the Beam-selected carrier schedule.
Its initializer is saved and reused byte-for-byte by the RGB control, so the resource
comparison changes supervision only.  Arm~1 carries \cref{claim:vram}.

\textbf{Arm 2 --- compact reconstruction with Beam Fusion.}
The same captures feed Beam Fusion (\cref{sec:beam-fusion}) and carrier refinement
(\cref{sec:carrier}), followed by the same compact optimizer under matched budgets.  Arm~2
carries only \cref{claim:beam}.

\textbf{RGB control.}
The repository-standard image-supervised 3DGS trainer, loading the exact frozen initialization
produced by Arm~1, with standard densification under a hard cap, on the same device,
resolution, seeds, and stopping rule.  The control may render and score the immutable compact
outputs; the compact processes never see its images.  This control is \emph{not} Original 3DGS:
Original 3DGS requires a real COLMAP sparse reconstruction for initialization, and a random or
Beam-derived initialization must not be relabeled as SfM.
\toshow{Original 3DGS and Original 3DGS with compressed RGB enter as separate baselines once a
real COLMAP model exists for each evaluation scene (obligation O-7, \cref{sec:obligations}).}

\subsection{Predeclared decision rule}
\label{sec:decision-rule}

\begin{itemize}[leftmargin=1.6em,itemsep=1pt]
  \item If Arm 1 reaches the preregistered held-out parity floor and materially lowers peak
        VRAM, the systems claim proceeds regardless of Beam.
  \item If Arm 2 also beats Arm 1 under the matched Beam gate (\cref{sec:exp2-decision}), Beam
        becomes an additional method contribution.
  \item If Arm 2 ties or loses, the direct compact pipeline is retained and Beam is reported as
        a bounded negative or exploratory result; it is not folded into the systems claim.
\end{itemize}

\begin{ToShowBlock}[Rule constants not yet frozen]
The parity floor (dB band on held-out PSNR plus SSIM/LPIPS guards), the materiality threshold
for ``materially lowers peak VRAM'', and the multiplicity policy across scenes and seeds must
be frozen in a preregistration \emph{before} the confirmatory runs, by a reviewer without
outcome access.  The draft blockers listed in
\repopath{docs/THREE_ARM_EXPERIMENT_PROGRAM.md} (exact compact optimizer, topology policy,
stopping rule for Arm 1) must be resolved in the same preregistration.
\end{ToShowBlock}
